\section{Establecimiento autónomo de objetivos y aprendizaje por currículo}
\subsection{Banco de objetivos y curriculum jerárquico}
Chelsea utiliza un «goal replay buffer» para registrar los objetivos alcanzados o casi alcanzados y, junto con una \textbf{curriculum policy}, selecciona para el entrenamiento objetivos de dificultad moderada. Cada vez que la policy alcanza un objetivo, este se escribe, junto con sus indicadores de capacidad, en el «tracker de estado de objetivos».

\begin{figure}[H]
\centering
\includegraphics[width=0.88\textwidth]{slides-images/slide-05.jpg}
\caption{La Slide 5 explica cómo implementar un curriculum automático en goal-conditioned RL.}
\end{figure}

\begin{tabularx}{\textwidth}{lX}
\toprule
Etapa & Medida principal \\
\midrule
Exploración & Al muestrear objetivos de entrenamiento, se da preferencia a novel/uncertain states; \\
Equilibrio & Se fija una success rate window (por ejemplo, 50\%-80\%) para mantener la zona de aprendizaje; \\
Avance & Por cada objetivo logrado, se aplica un parameterized difficulty increment; \\
\bottomrule
\end{tabularx}

\subsection{Goal-Query y el razonamiento del agente}
El ponente mencionó el módulo «Goal-Query»: durante el entrenamiento, el agent se pregunta continuamente «¿qué skill debo completar para poder generalizar después?» y elige objetivos con base en un novelty score. Este proceso cuida el equilibrio entre exploration (novel targets) y exploitation (objetivos conocidos).

\begin{knowledgebox}{Flujo de Goal-Query}
1. Se calcula la novelty de cada stored goal mediante una uncertainty estimate;\\
2. Si la novelty supera el threshold, se programa un exploration goal;\\
3. Con apoyo del skill graph se infieren las dependencias entre goals y se decide si conviene aprender primero los goals básicos.
\end{knowledgebox}

\subsection{Goal space engineering}
Chelsea enfatiza que el goal space debe admitir un escalamiento dinámico, con centro/radio/priority definidos para cada objetivo. Por ejemplo, un complex goal se descompone en primitives; durante el entrenamiento se muestrean varias veces goals de nivel Primitive y luego un scheduler avanza hacia el composite goal.

\begin{tabularx}{\textwidth}{lX}
\toprule
Estrategia & Función \\
\midrule
Primitive sampling & Genera reward-distinct goals a partir de combinaciones de acciones básicas; \\
Composite scheduler & Cuando la primitive success rate alcanza la meta, pondera proporcionalmente la probabilidad de selección del composite goal; \\
Priority annealing & Ajusta dinámicamente la goal priority según la novelty/entropy; \\
\bottomrule
\end{tabularx}

\subsection{Resumen de la sección}
El establecimiento autónomo de objetivos combina el aprendizaje por currículo, Goal-Query y el difficulty scheduler, de modo que el robot puede acumular habilidades de forma gradual sin que una persona le asigne tareas.
\newpage
